\chapter{Architecture}
\label{ch:architecture}

\section{The shape of the application}

The application is a server-rendered web application. There is no separate
front-end project and no JSON API between the pages and the database. A browser
asks for a page, FastAPI runs a handler, the handler calls a service, and the
service returns plain data that a Jinja template renders.

\begin{figure}[h]
\centering
\resizebox{\textwidth}{!}{%
\begin{tikzpicture}[
  node distance=8mm,
  box/.style={draw=accent!60, fill=accent!5, rounded corners=2pt,
              minimum height=9mm, minimum width=26mm, align=center, font=\small},
  wide/.style={box, minimum width=38mm},
  arrow/.style={-{Latex[length=2mm]}, accent2, thick}
]
\node[box] (browser) {Browser};
\node[box, below=of browser] (tg) {Telegram};
\node[wide, right=18mm of browser] (routers) {FastAPI routers\\\scriptsize\code{app/routers/}};
\node[wide, right=18mm of tg] (bot) {Bot dispatcher\\\scriptsize\code{services/bot.py}};
\node[wide, right=16mm of routers, yshift=-9mm] (services) {Services\\\scriptsize parsers, pnl, aggregation,\\\scriptsize prices, reports, backup};
\node[box, right=16mm of services] (orm) {SQLAlchemy};
\node[box, below=of orm] (db) {Postgres\\\scriptsize or SQLite};

\draw[arrow] (browser) -- (routers);
\draw[arrow] (tg) -- (bot);
\draw[arrow] (routers) -- (services);
\draw[arrow] (bot) -- (services);
\draw[arrow] (services) -- (orm);
\draw[arrow] (orm) -- (db);
\end{tikzpicture}}
\caption{Both entry points, the web pages and the bot, share the same services.}
\end{figure}

The important line in that diagram is the one from the bot into the services.
The bot is not a second implementation of the application; it is a second
interface to the same functions. When a total changes, it changes in one place
and every interface follows.

\section{Directory layout}

\begin{lstlisting}[style=sh]
app/
  main.py         application object, middleware, lifespan, health
  config.py       settings read from the environment
  db.py           engine and session factory
  models.py       ORM models
  auth.py         login, logout, route guard
  security.py     password hashing, CSRF, login rate limit
  templating.py   shared Jinja instance and filters
  format.py       money and percentage formatting
  scheduler.py    in-process scheduled jobs
  routers/        one module per section of the site
  services/       every rule that decides a number
  templates/      Jinja pages
  static/         compiled CSS
migrations/       Alembic
scripts/          seed.py, backup.py
tests/
\end{lstlisting}

\section{The rule that keeps it tidy}

Route handlers do four things and nothing else: read the form, check the CSRF
token, call a service, and redirect or render. They never contain arithmetic.

The reason is not style. The Telegram bot and the PDF generator do not go
through routers at all, so any calculation written inside a handler would be
invisible to them and would eventually be duplicated. Duplicated arithmetic on
money drifts, and drifting numbers are hard to notice and harder to trust.

\begin{keypoint}[Where logic lives]
If a piece of code decides an amount, a total, a percentage or a category, it
belongs in \file{app/services/}. If it decides a redirect or an HTTP status, it
belongs in \file{app/routers/}.
\end{keypoint}

\section{The service modules}

\begin{table}[h]
\centering
\begin{tabular}{@{}p{4.2cm}p{9.5cm}@{}}
\toprule
\textbf{Module} & \textbf{Responsibility} \\
\midrule
\code{parsers/} & Turn a broker file into a list of positions. One module per
broker, plus a dispatcher. \\
\code{pnl.py} & Profit of a position. The formula depends on the broker. \\
\code{aggregation.py} & Every total the application shows: KPIs, chart series,
net worth snapshots. \\
\code{prices.py} & Market prices, and the mapping from a crypto pair to a
ticker. \\
\code{symbol\_lookup.py} & From an ISIN to a market symbol, with a cache. \\
\code{contributions.py} & One-off and monthly contributions. \\
\code{expenses\_parse.py} & The paste parser and its safe arithmetic. \\
\code{reports/} & Excel, PDF and the charts embedded in them. \\
\code{bot.py} & Telegram: dispatch, replies, keyboards, confirmation. \\
\code{backup.py} & Database dump and delivery. \\
\bottomrule
\end{tabular}
\caption{The service layer.}
\end{table}

\section{Configuration}

All configuration is read once, at import time, into a single settings object
built with \code{pydantic-settings}:

\begin{lstlisting}[style=py]
class Settings(BaseSettings):
    model_config = SettingsConfigDict(env_file=APP_ROOT / ".env", extra="ignore")

    database_url: str = ""          # empty -> local SQLite
    secret_key: str = "dev-insecure-change-me"
    cookie_secure: bool = False
    is_prod: bool = False
    report_owner: str = "Account owner"
    currency_symbol: str = "€"
    ...

    @property
    def sqlalchemy_url(self) -> str:
        if self.database_url.strip():
            return self.database_url.strip()
        return f"sqlite:///{APP_ROOT / 'dev.db'}"
\end{lstlisting}

Two properties matter more than the rest. \code{sqlalchemy\_url} makes SQLite
the default, so a fresh clone runs with no database to install. And
\code{telegram\_enabled} is simply whether a token was provided, which is what
decides at boot whether the bot router is mounted at all.

\begin{keypoint}[Optional features are optional]
The bot, the scheduler and the task endpoints are all disabled by an empty
value: no token, no bot; no \code{TASKS\_TOKEN}, no task endpoints. A user who
only wants the web pages never has to switch anything off.
\end{keypoint}

\section{The application lifespan}

FastAPI runs a lifespan function on startup and shutdown. Three things happen
there:

\begin{lstlisting}[style=py]
@asynccontextmanager
async def lifespan(app: FastAPI):
    if settings.is_prod and settings.secret_key == "dev-insecure-change-me":
        raise RuntimeError("SECRET_KEY is not configured in production.")
    scheduler = None
    if settings.use_scheduler:
        from app.scheduler import start_scheduler
        scheduler = start_scheduler()
    if settings.telegram_enabled:
        from app.services.bot import setup_webhook
        await setup_webhook()
    yield
    if scheduler:
        scheduler.shutdown(wait=False)
\end{lstlisting}

The first check refuses to boot a production instance with the default signing
key. A known key means any visitor can forge a session cookie, so failing to
start is the correct response, not a warning in a log nobody reads.

The webhook registration is repeated on every boot on purpose. It is idempotent,
and it means a redeploy is enough to recover a bot that lost its webhook.

\section{Two health endpoints, on purpose}

The application exposes \code{/health} and \code{/health/db}, and they are not
interchangeable.

\begin{itemize}
  \item \code{/health} answers always and never touches the database. It is the
        liveness check for the hosting platform.
  \item \code{/health/db} runs \code{SELECT 1} and answers \code{503} when the
        database does not reply. It is the readiness check for an external
        monitor.
\end{itemize}

\begin{pitfall}[Why not one endpoint]
If the platform healthcheck queried the database, a brief database outage would
make the platform consider the application unhealthy and tear the deployment
down. The application would then be gone even after the database came back. The
split keeps a database problem a database problem.
\end{pitfall}
